\scanpage{5}

7) Modules pour $S\subset K\subset X$

On suppose \inserted{\deleted{$X$} compacte, et} pour simplifier $X$ \deleted{compacte} \inserted{sans bord} orientée. Ceci dit, on fixe $S_0\subset K_0\subset X_0$ situation top, et on \uline{rigidifie} $(S\subset K\subset X)$ en se donnant une classe d'isotopie d'\unsure{hom.|isom.|homéo.} top. avec la situation standart. On élimine ainsi les autos, sauf dans le cas \deleted{$\unsure{X_0|K_0|X}$} où $X_0$ a des composantes de degré $0,1$.....

Pour avoir un ``\unsure{nb|ens} fini de modules'', on \inserted{pourrait} exiger \inserted{p. ex.} que les arêtes de $K$ soient des arcs de cercles pour la conn. projective canonique. Mais les ``cartes \inserted{cell.} holomorphes'' ne satisfont pas néc. à cette condition ! On pourrait exiger plutôt que pour \deleted{\unsure{une|ces|avec}} les : paramétrisations géodésiques sur la découpe, sur disques \unsure{que|orc|avec} l'\unsure{homéomorphisme|homographie|isomorphisme} de recollement soit linéaire. [Éventuelt on pourrait exiger l'équidistance \inserted{des sommets} sur chaque composante connexe du bord de la découpe...] Mais on a des ennui pour les morceaux \deleted{de genre 0} \inserted{qui sont des disques \unsure{non ponctués|monoponctués}} il n'y a pas de métrique \unsure{inst courb $-$|intrinsèque|constant}. Il

% archive sheet number: 4
